\documentclass[11pt,a4paper]{article}

\usepackage[utf8]{inputenc}
\usepackage[T1]{fontenc}
\usepackage[english]{babel}
\usepackage[margin=2.0cm]{geometry}
\usepackage{parskip}
\usepackage{hyperref}
\usepackage{csquotes}
\usepackage[backend=biber,style=numeric-comp,sorting=none,maxbibnames=99]{biblatex}

\addbibresource{refs.bib}
\hypersetup{colorlinks=true,urlcolor=blue,linkcolor=black}

\setlength{\parskip}{0.45em}
\AtBeginBibliography{\footnotesize}
\setlength{\bibitemsep}{0.05em}

\title{\vspace{-1.4cm}Pre-assignment: AI Solution Sprint 2026}
\author{Tauheed Butt \\ \small \href{https://tauheed.dev}{tauheed.dev}}
\date{6 September 2026}

\begin{document}

\maketitle

\section*{What is a hackathon?}

A hackathon is a short event where people form teams and build something in a
fixed amount of time. The limit is usually one to three days. The teams are mixed
on purpose, so a developer, a designer and a business student end up working on
the same problem. The problem normally comes from a real company, not from a
textbook. At the end each team shows what it built \cite{nolte2023}.

The way of working is different from a course project. There is no time to plan
everything first. You pick a small piece, build it, show it, and fix it.
Decisions get made quickly because the clock is always visible. Teams also talk
to mentors from the partner company, so feedback comes back fast. Research on
hackathons in education reports the same effects. People learn by doing, and the
mixed teams are a large part of why \cite{medina2020}.

I have seen this myself. Last year I took part in the Since AI Hackathon in
Turku. It ran for 72 hours with more than 260 builders, and the rule was simple:
no slide decks, only working demos \cite{sinceai_projects}. My team won the
ElevenLabs challenge. What I learned there was that the hardest part is not the
code. It is deciding fast what not to build.

\section*{My goals for the AI Solution Sprint}

My first goal is to build something the partner company can actually use. Not a
demo that stops working on Thursday evening, but a working piece of software they
could take into further development.

My second goal follows from the first. I would like to keep working with that
company after the event. Two days is a good way to start something longer, and
that is what I am hoping for.

My third goal is about the team. I have mostly built things alone. I want to see
how much faster a mixed team can move when the people in it are not all
developers.

\section*{Benefits for everyone involved}

\textbf{For participants.} You get exposure and real practice. You build
something under pressure with people you have never met, which is close to how
real work feels. You also meet companies directly, and that can turn into
something longer term.

\textbf{For the challenge owners.} They get fresh ideas on a problem they are
already stuck on. Several teams work on the same challenge at once, so they see
several directions in two days. Sometimes one of them is good enough to take into
development. They also meet people they might later hire.

\textbf{For everyone else.} The universities get students working on real
industry problems. The cities behind AIStart get a stronger local tech scene
\cite{aistart}. The venue and the sponsors get visibility, and students new
to the field get to see how the tech industry actually works.

\section*{The challenge that interests me most}

I would like to work on the Rollyy challenge. Rollyy is a Finnish startup that
builds mobile charging robots. You park your car, book a session, and the robot
drives to your car and charges it \cite{rollyy}.

Their challenge is about a bigger problem. Charging an electric car today means
collecting apps, because every charging company wants its own app and account.
Rollyy wants one pass instead. You book, you get access, and you pay, all in one
flow. The hard part is that this pass must not become one more app.

My idea is to keep it out of any new app. Booking would happen through an MCP
server and a voice agent, so the driver just asks for a charger. Access would
work through a pass in Apple Wallet or Google Wallet that the driver taps at the
charger \cite{applepasskit}. Payment would use a card saved once through
Apple Pay or Google Pay, charged after the session ends. Behind all of it, Rollyy
would act as the service provider holding the driver's account, which is how
charging networks already share customers with each other \cite{ocpi_guide}.

The parts I feel confident about are the MCP server, the agents, the payment
integration and the cloud setup. I have built all of those before. The part I
find hard is access. I have not generated signed wallet passes before, and I do
not yet know what identity Rollyy's own API expects. That is where I would want
mentor time on the first day.

My second choice is the Cognovo challenge, because finding and ranking manual
work inside a company is a clean problem for AI agents.

\section*{What I can bring to a team}

I have over four years of full stack development experience, mostly TypeScript,
React and Node, with PostgreSQL and AWS behind it \cite{tauheeddev}. Lately
I have worked on AI agents: building them with AI SDKs, writing MCP servers, and
using embeddings and retrieval. I also do product design, so I can decide what a
screen should do, not only how to build it.

In a mixed team my most useful skill is translating. I can listen to a business
idea and say quickly whether it can be built in two days, and what a smaller
version of it would look like. In a short event that saves a lot of time.

\section*{How I used AI in this assignment}

I used an AI assistant as a research and drafting tool. I gave it my own goals,
my experience and my own view on the challenges, and asked it to find sources and
help me structure the report.

I used it most on the Rollyy challenge. The brief confused me at first, so I went
back and forth with the assistant, arguing my own reading against it until the
challenge became clear. That is how I reached the plan above. I then wrote and
edited the text myself, and I opened and checked every source. I did not use AI
output as a source.

\printbibliography[title={Sources}]

\end{document}
